\mainlanguage[de]
\input include/conventions/units.tex
\input include/workarounds/lmtx-units-2026.tex
\starttext
\chapter{Einheiten -- Regressionstests}

\section{Text und Mathematik}
Im Messbereich \m{\unit{2 kilo ohm}} ergibt sich
\m{\unit{0.152 kilo ohm}}.

\m{R = \unit{5.6 kilo ohm}},
\m{f(1,2)} und
\m{U = \unit{100 micro volt}}.

\section{Ziffern und Vorzeichen}
\m{\unit{3.1415}},
\m{\unit{-0.001}},
\m{\unit{1.2e-3 ampere}},
\m{\unit{0.5 percent}}.

\section{Zusammengesetzte Groessen}
\m{\unit{4.25 euro}},
\m{\unit{0.5 volt per division}},
\m{\unit{1 kilogram square meter per ampere cubic second}}.

\startformula
 \frac{\unit{1 meter}-\unit{3 meter}}
      {\unit{4 second}-\unit{6 second}}
 \qquad
 \unit{1 volt}
 = \frac{\unit{kilogram}\cdot\unit{square meter}}
        {\unit{ampere}\cdot\unit{cubic second}}
\stopformula

\section{MetaPost (feste TeX-Beschriftung)}
\startMPcode
 draw textext(btex \m{\unit{5.6 kilo ohm}} etex);
\stopMPcode
\stoptext
